\chapter{배선표와 스크램블러: \texorpdfstring{\code{wirings\_pkg}, \code{scrambler}}{wirings\_pkg, scrambler}}\label{ch:rtl-scrambler}

G부는 B부의 모델을 회로로 옮긴 기준 RTL(\code{rtl/src})을 읽는다. 파일은 네 개다.

\begin{center}
\small
\begin{tabular}{lll}
\toprule
파일 & 하는 일 & Rust에서는 \\
\midrule
\code{wirings\_pkg.sv} & 배선표와 로터 통과 함수 & \code{wiring}, \code{through} \\
\code{scrambler.sv} & 드럼 한 벌의 치환표 (조합 논리) & \code{WheelOrder::scrambler} \\
\code{bombe\_core.sv} & 한 드럼 위치의 판정 & \code{Plugged::propagate} \\
\code{bombe\_sweep.sv} & 모든 위치를 돌며 정지를 내보냄 & \code{run\_slice} \\
\bottomrule
\end{tabular}
\end{center}

기준 RTL은 시뮬레이션용이다(\ref{sec:sim-vs-synth}절). 칩에 올리는 것이 아니라 ``클럭 단위로 본 올바른
동작''을 정의하는 코드이므로, 크기보다 읽기 쉬움과 Rust 모델과의 대응을 우선했다.

\section{약속: 글자, 번호, 비트의 배치}

회로에는 ``글자''라는 자료형이 없다. 모든 것이 비트의 묶음이다. 그래서 먼저 약속을 정한다.

\begin{itemize}
\item 글자는 5비트 값 0--25다(\texttt{A}=0). 5비트는 0--31을 나타낼 수 있으므로 26--31은 쓰이지 않는 값이다.
\item 로터는 번호로 가리킨다(0=I, \ldots, 7=VIII, 8=Beta, 9=Gamma). 반사판도 번호다(0=B, 1=C, 2=얇은 B, 3=얇은 C).
\item 676비트 상태에서 케이블 $c$의 $w$번 선은 $26c + w$번째 비트다.
\end{itemize}

\Why 약속을 문서로 못박아 두는 것은 검증 때문이다. 테스트벤치가 Rust 모델의 \code{live[c]}의 $w$번째 비트와
회로의 어느 비트를 비교해야 하는지가 이 약속으로 정해진다. 약속이 어긋나면 양쪽이 다 옳아도 비교는 실패한다.

\section{배선표: 생성된 상수}

\begin{lstlisting}[language=SV,caption={\code{wirings\_pkg.sv}의 머리(상수의 값은 생략).}]
// GENERATED by rtl/gen/gen_wirings.py from data/wirings/enigma.txt -- do not edit.
// Tables are packed: entry [id][c] lives at bits ((id*26 + c)*5) +: 5.
package wirings_pkg;
  localparam int NUM_ROTORS = 10;
  localparam int NUM_REFLECTORS = 4;
  localparam logic [NUM_ROTORS*26*5-1:0] ROTOR_FWD = 1300'h...;
  localparam logic [NUM_ROTORS*26*5-1:0] ROTOR_BWD = 1300'h...;
  localparam logic [NUM_REFLECTORS*26*5-1:0] REFLECTOR = 520'h...;
\end{lstlisting}

\code{package}는 여러 모듈이 함께 쓰는 상수와 함수를 담는 곳이다. 로터 10개 × 글자 26개 × 5비트 = 1,300비트의
상수 하나에 모든 로터의 정방향 배선표가 차곡차곡 들어 있다. \code{x[i +: 5]}는 ``$i$번째 비트부터 5비트''라는
뜻이다. 로터 \code{id}의 글자 \code{c}에 대한 항목은 \code{(id*26 + c)*5}번째 비트부터 5비트다.

첫 줄의 주석이 중요하다. 이 파일은 사람이 쓴 것이 아니라 스크립트가 \code{data/wirings/enigma.txt}에서
만든 것이다. Rust 모델이 읽는 바로 그 파일이다.

\General 회로에서 바뀌지 않는 표는 이렇게 상수로 적는다. 합성하면 표의 내용이 LUT에 굳어진다. 7세그먼트
표시기에 숫자를 그리는 회로(숫자 0--9마다 어느 획을 켤지의 표), 글꼴 표, 사인 표가 모두 같은 방식이다.

\section{26으로 나눈 나머지를 나눗셈 없이}

\begin{lstlisting}[language=SV,caption={\code{wirings\_pkg.sv}: 26진 덧셈과 뺄셈.},label={lst:add26}]
// (a + b) mod 26 for letters a, b in 0..25.
function automatic [4:0] add26(input [4:0] a, input [4:0] b);
  reg [5:0] s;
  begin
    s = {1'b0, a} + {1'b0, b};
    if (s >= 6'd26) s = s - 6'd26;
    add26 = s[4:0];
  end
endfunction

// (a - b) mod 26 for letters a, b in 0..25.
function automatic [4:0] sub26(input [4:0] a, input [4:0] b);
  reg [5:0] s;
  begin
    s = {1'b0, a} + 6'd26 - {1'b0, b};
    if (s >= 6'd26) s = s - 6'd26;
    sub26 = s[4:0];
  end
endfunction
\end{lstlisting}

\code{\{1'b0, a\}}는 5비트 \code{a} 앞에 0을 한 비트 붙여 6비트로 만든 것이다. 두 글자의 합은 최대 50이라
5비트(최대 31)에 담기지 않기 때문이다.

\Why Rust에서는 \code{(c + s) \% 26}이라고 썼다. 회로에서 \code{\%}는 나눗셈 회로를 뜻하고, 나눗셈은 덧셈보다
훨씬 크고 느리다. 그런데 여기서는 나눗셈이 필요 없다. 두 수가 모두 0--25이면 합은 0--50이고, 26으로 나눈
나머지는 ``26 이상이면 26을 한 번 뺀다''로 끝난다. 비교 한 번, 뺄셈 한 번, 선택 한 번이다. 입력의 범위를
알고 있으면 일반적인 연산을 훨씬 싼 회로로 바꿀 수 있다.

\General 시계의 시침이 같다. 9시에서 5시간 뒤는 14시가 아니라 2시다. 12를 넘으면 12를 한 번 빼면 된다.
14를 12로 ``나눌'' 필요가 없다.

\Ours 같은 파일의 주석에 이 함수들이 조금 옛 문법으로 쓰인 이유가 적혀 있다. 시뮬레이터 두 종과 합성 도구
한 종이 모두 받아들이게 하기 위해서다. 도구마다 지원하는 문법이 조금씩 다르므로, 여러 도구에서 같은 결과를
확인하려면 가장 보수적인 문법을 쓴다.

\section{로터 한 단: \texorpdfstring{\code{rotor\_fwd}}{rotor\_fwd}}

\begin{lstlisting}[language=SV,caption={\code{wirings\_pkg.sv}: 회전한 로터를 통과한다.},label={lst:rotorfwd}]
function automatic [4:0] rotor_fwd(input [3:0] id, input [4:0] core, input [4:0] c);
  reg [4:0] x;
  begin
    x = ROTOR_FWD[({6'd0, id} * 10'd26 + {5'd0, add26(c, core)}) * 5 +: 5];
    rotor_fwd = sub26(x, core);
  end
endfunction
\end{lstlisting}

코드~\ref{lst:through}의 \code{through}와 나란히 놓고 읽으면 같은 세 토막이다. \code{add26(c, core)}로
밀고, 표에서 읽고, \code{sub26(x, core)}로 되민다.

달라진 것은 ``표에서 읽는다''의 뜻이다. Rust에서는 메모리의 한 칸을 읽는 것이었다. 회로에서는
\ref{sec:mux-cost}절의 선택 회로다. 로터 번호 \code{id}와 글자가 실행 중에 정해지므로, 이 한 줄은 260개 항목
(로터 10개 × 글자 26개) 가운데 하나를 고르는 5비트 폭의 선택 회로가 된다(그림~\ref{fig:g-rotorfwd}).

\begin{figure}[ht]
\centering
\begin{tikzpicture}[x=1cm,y=1cm]
  \node[font=\ttfamily\small] (c) at (0,0.4) {c};
  \node[font=\ttfamily\small] (core) at (0,-0.6) {core};
  \node[blkb,minimum width=18mm] (add) at (2.2,0) {add26};
  \node[mux,minimum width=22mm,minimum height=11mm] (rom) at (5.6,0) {260:1};
  \node[note,above=1mm of rom] {배선표 (상수)};
  \node[blkb,minimum width=18mm] (sub) at (9.0,0) {sub26};
  \node[font=\ttfamily\small] (o) at (11.4,0) {rotor\_fwd};
  \draw[->] (c) -- ([yshift=2mm]add.west); \draw[->] (core) -- ([yshift=-2mm]add.west);
  \draw[->] (add) -- (rom); \draw[->] (rom) -- (sub); \draw[->] (sub) -- (o);
  \node[font=\ttfamily\small] (id) at (5.6,-1.7) {id};
  \draw[->] (id) -- (rom.south);
  \draw[->] (core.east) -- ++(0.3,0) |- (9.0,-1.2) -- (sub.south);
\end{tikzpicture}
\caption{\code{rotor\_fwd}가 만드는 회로. 5비트 폭의 선이 지난다.}
\label{fig:g-rotorfwd}
\end{figure}

\section{스크램블러}

코드~\ref{lst:g-scr}(\ref{ch:hdl}장)에서 이미 본 모듈이다. 입출력만 다시 본다.

\begin{lstlisting}[language=SV,caption={\code{scrambler.sv}의 입출력.}]
module scrambler (
  input  logic [1:0]       refl_id,
  input  logic             greek_en,
  input  logic [3:0]       greek_id, left_id, middle_id, right_id,
  input  logic [4:0]       core_g, core_l, core_m, core_r,
  output logic [25:0][4:0] perm
);
\end{lstlisting}

\code{module}은 부품의 한 종류를 정의한다. \code{input}과 \code{output}은 부품의 단자다. 입력은 로터 순서
(반사판과 로터의 번호)와 네 로터의 코어 위치이고, 출력 \code{perm}은 5비트 값 26개, 곧 26글자 각각이 어느
글자로 바뀌는지다. 클럭 입력이 없다. 순수한 조합 논리여서 입력이 바뀌면 출력이 따라 바뀔 뿐이다.
\code{greek\_en}은 M4의 네 번째 로터를 쓸지를 정한다. 꺼져 있으면 그리스 로터의 통과가 건너뛰어진다.

이 모듈 하나가 얼마나 큰지 세어 보자. 글자 26개마다 로터 통과가 여섯 번에서 여덟 번이고, 통과 한 번이
260:1 선택 회로(5비트)에 덧셈과 뺄셈이다. 합성하면 LUT 19,717개다. \code{bombe\_core}는 이것을 메뉴의
선마다 하나씩 둔다.

\Ours 실제 봄베와 견주면 이 모듈은 드럼 세 개 한 벌이다. 다만 실제 드럼은 전선으로 되어 있어 26가닥의
전류를 양방향으로 그대로 통과시키지만, 이 모듈은 ``어느 글자가 어느 글자로 가는가''라는 \emph{표}를 내놓는다.
전류를 실제로 옮기는 일은 다음 장의 전파 논리가 이 표를 보고 한다.

\paragraph{이 모듈의 시험.} 테스트벤치는 무작위로 고른 로터 순서와 위치 3,000건에 대해 26글자의 출력을
Python으로 쓴 참조 계산과 비교한다. 여기에 오른쪽 로터의 모든 위치, 그리고 \texttt{AAAAA}가 \texttt{BDZGO}가
되는 잘 알려진 값을 더했다.
